\section{Method: from atlas point to falsifiable experiment}
\label{sec:method}

\subsection{PhaseCard: one typed experimental record}

The machine-readable PhaseCard joins the scientific tuple in
Eq.~\eqref{eq:atlas-tuple} to a protocol and an epistemic record:
\begin{equation}
  \mathcal P=
  (\underbrace{D,X,\delta,\Omega,\mathbf N,m,H}_{\text{scientific object}},
   \underbrace{g,\mathcal I,\mathcal N,\mathcal H,\mathcal F}_{\text{protocol}},
   \underbrace{T,L,\mathbf E,C}_{\text{epistemic record}}).
  \label{eq:phasecard}
\end{equation}
Here $g$ denotes controls, $\mathcal I$ interventions, $\mathcal N$ nulls,
$\mathcal H$ holdouts, and $\mathcal F$ falsifiers. $T$ is theory status, $L$
fidelity, $\mathbf E$ a vector of evidence coordinates, and $C$ the outcome.
Resources, software identity, seed streams, and artifact hashes are attached as
execution metadata; changing them does not silently change the scientific
identity.

This representation enforces several invariants. Inner stochastic replicas are
not counted as independent outer disorder. A threshold has an explicit crossing
or censor status. A metric definition includes units and normalization. A
post-hoc selected coordinate cannot be labeled preregistered. A failed
accelerator attempt and a fallback are distinct execution attempts under one
scientific condition.

\subsection{An observable is more than a scalar}

Every reported estimate has the form
\begin{equation}
  \widehat m=(\bar m,[m_-,m_+],q,U,n_{\mathrm{outer}},
  n_{\mathrm{inner}},\Omega,\text{units},\text{status}),
  \label{eq:typed-estimate}
\end{equation}
where $q$ is interval level and $U$ names the uncertainty procedure. The status
records validity and censoring. This prevents a temporal standard deviation, a
between-initialization variance, and a bootstrap interval from sharing the same
unqualified field name.

\manuscriptfigure{figures/figure3_observable_map.pdf}{Observable dictionary used
throughout the atlas. Risk, semantic recovery, specialization, response,
fluctuation, dynamics, and intervention effects require different averaging
ensembles and nulls. Similar numerical ranges do not make the estimands
interchangeable.}{fig:observable-map}

\subsection{Primary observables}

\paragraph{Intensive predictive risks.}
Cross entropy is normalized by the random-reference entropy,
$R_{\mathrm{CE}}=\ell_{\mathrm{CE}}/\log V$, and Brier risk by $1-1/V$ for
vocabulary size $V$. Train, IID, and OOD risks are stored separately:
\begin{equation}
  R_{\mathrm{train}},\qquad R_{\mathrm{IID}},\qquad R_{\mathrm{OOD}},
  \qquad e_{\mathrm{gen}}=R_{\mathrm{IID}}-R_{\mathrm{train}}.
\end{equation}
Tokens, FLOPs, wall time, parameter count, and examples are resources rather
than extensive versions of the risk.

\paragraph{Semantic order.}
When a latent variable $Z$ is known, semantic retention is measured through a
registered probe, conceptually
\begin{equation}
  S_{\mathrm{sem}}=I(m;Z)/H(Z),
\end{equation}
and checked against shuffled labels. In the public reference grid the stored
semantic order is the signed task-success coordinate produced by the controlled
retrieval task. It remains a task order, not a universal magnetization.

\paragraph{Outer-disorder response and Binder diagnostic.}
For independent outer-seed terminal orders $m_s$ and declared finite-size
coordinate $N$,
\begin{equation}
  \chi_N=N\left(\langle m^2\rangle_{\Omega}
  -\langle m\rangle_{\Omega}^2\right),\qquad
  U_4=1-\frac{\langle m^4\rangle_{\Omega}}
  {3\langle m^2\rangle_{\Omega}^2}.
  \label{eq:phase-diagnostics}
\end{equation}
Checkpoint fluctuations from one training trajectory are named temporal
variance and never substituted into Eq.~\eqref{eq:phase-diagnostics}.

\paragraph{Causal module effect.}
For matched full and ablated evaluations, the raw effect is
$\Delta R=R_{\mathrm{ablated}}-R_{\mathrm{full}}$. A bounded signed effect uses a
non-vanishing registered risk scale. Ratios that divide by remaining accuracy
near zero are invalid because they can diverge without a physical change.

\subsection{Outcome assignment and finite-size logic}

A phase-language claim must pass five gates: semantic validity; an independent
disorder ensemble; an interior finite-size signature such as a crossing,
sharpening response, or stable collapse; robustness to fitting window and
minimum size; and a frozen prospective check. Failing a gate produces a named
negative outcome rather than an omitted figure.

\manuscriptfigure{figures/figure4_phase_decision.pdf}{Decision rule for
finite-size claims. Each failed gate has a specific negative outcome. Passing
the sequence licenses only a domain-specific, bounded claim; it does not by
itself establish universality or an equilibrium thermodynamic limit.}
{fig:phase-decision}

\subsection{Predictive phase continuation}

Predictive continuation is one protocol inside the atlas. Given calibration
observables $\mathcal O_{\mathrm{cal}}$ and a surrogate
$\mathcal O_{\mathrm{solv}}(\theta)$, fit
\begin{equation}
  \widehat\theta_{\mathrm{eff}}=arg\min_{\theta}
  D\!\left(\mathcal O_{\mathrm{cal}}^{\mathrm{target}},
  \mathcal O_{\mathrm{solv}}(\theta)\right),
\end{equation}
and evaluate only on withheld quantities,
\begin{equation}
  \Ebridge=D\!\left(\mathcal O_{\mathrm{holdout}}^{\mathrm{target}},
  \mathcal O_{\mathrm{solv}}(\widehat\theta_{\mathrm{eff}})\right).
\end{equation}
Constant, source-boundary, size-only, nearest-calibration, and base-surrogate
predictors are mandatory alternatives. Crossed, left-censored, and
right-censored thresholds are retained separately. The controlled simulator in
this release tests this accounting logic; it does not establish transport in a
real training stack.

\FloatBarrier
